% =====================================================
% Introduction
% =====================================================
\section{Introduction}
Diabetic Retinopathy (DR) is a complication of diabetes and a leading cause of preventable blindness. It is graded on a five-point scale, from no retinopathy (Grade 0), through mild, moderate, and severe non-proliferative DR, to proliferative DR (Grade 4). Early stages are asymptomatic, so screening relies on colour fundus photography, where clinicians look for microaneurysms, haemorrhages, exudates, and neovascularisation. Manual grading is slow, requires specialist expertise, and varies between graders, which has motivated automated deep-learning systems that grade fundus images directly.

For this work we use the publicly available \textbf{APTOS 2019 Blindness Detection} dataset (Kaggle), containing $3{,}662$ high-resolution colour fundus photographs from rural India, labelled by clinicians on the five-class International Clinical Diabetic Retinopathy (ICDR) severity scale ($0$--$4$).
